\documentclass[final, hyperref, table]{beamer}
\mode<presentation>


 %\usepackage[english]{babel} % "babel.sty"
% \usepackage{french}                  % "french.sty"
%  \usepackage{franglais}               % "franglais.sty" (a defaut)
  \usepackage{times}			% ajout times le 30 mai 2003
 
%% --------------------------------------------------------------
%% CODAGE DE POLICES ?
%% Si votre moteur Latex est francise, il est conseille
%% d'utiliser le codage de police T1 pour faciliter la césure,
%% si vous disposez de ces polices (DC/EC)
\usepackage[utf8]{inputenc}
\usepackage[T1]{fontenc}


%% ==============================================================
%\usepackage{graphicx}
\usepackage{amsmath,amsfonts}
%\usepackage[table]{xcolor}
\usepackage{subfigure}
\usepackage{fancybox}
%\usepackage{hyperref}
\usepackage{multicol}
\usepackage{wrapfig}
\usepackage{listings}
\usepackage{xcolor}

\usetheme{Warsaw}
%\usetheme{JuanLesPins}
\setbeamercovered{transparent}


% telemeta red
\definecolor{telemetaRed}{rgb}{0.41568, 0.01176, 0.02745}	% #6A0307
\usecolortheme[rgb={0.41568, 0.01176, 0.02745}]{structure} 

% Display a grid to help align images
%\beamertemplategridbackground[1cm]

%We will get the normal bibliography style (number or text instead of icon) by including the following code
\setbeamertemplate{bibliography item}[text]
\setbeamerfont{caption}{size=\footnotesize}
% listings settings
\definecolor{lstComments}{rgb}{0,0.6,0}
\definecolor{lstBkgrd}{rgb}{1,1,0.8}
\lstset{%
  language=Python, % the language of the code
  frame=single,  % adds a frame around the code
  commentstyle=\color{lstComments},% comment style
  backgroundcolor=\color{lstBkgrd},   % choose the background color
  basicstyle=\scriptsize,       % the size of the fonts that are used for the code
  keywordstyle=\color{blue},      % keyword style
  showstringspaces=false,          % underline spaces within strings only
}
\title[TELEMETA, audio web CMS for Ethnomusicological archives]{TELEMETA\\ Audio web Content Management System \\for ethnomusicological sound archives}

\author[Fillon, Pellerin, Brossier, Simonnot]{Thomas Fillon \inst{1,2}, Guillaume Pellerin\inst{1}, Paul Brossier\inst{1}, Jos{\'e}phine Simonnot\inst{3}}


\institute[Parisson]{\tiny
  \inst{1}%
  Parisson, Paris, France\\
  \inst{2}%
  LAM, Institut Jean Le Rond d'Alembert, UPMC Univ. Paris 06, UMR CNRS 7190, Paris, France\\
 \inst{3}%
  CREM, LESC, UMR CNRS 7186, MAE, Université Paris Ouest Nanterre La Défense, Nanterre, France\\
%Thanks
\vskip1ex
 {\tiny \textcolor{red}{\emph{This work was partially done inside the DIADEMS project\\ funded by the French National Research Agency ANR (CONTINT)}}}
 \begin{center}
\hfill
   \raisebox{-4ex}{\includegraphics[width=0.1\linewidth]{../poster/img/logo_CREM.png}} \hfill
  \includegraphics[width=0.15\linewidth]{../poster/img/logo_LESC.png}\hfill
   \includegraphics[width=.3\linewidth]{../poster/img/parisson_logo_FINALE_com.pdf}\hfill
   \includegraphics[width=.18\linewidth]{../poster/img/upmc.png}\hfill
 \end{center}

}
        

\begin{document}
\begin{frame}
  \maketitle
\end{frame}

\begin{frame}\tiny
\frametitle{Telemeta}
% ==================================
% --------- Résumé -----------------
% ==================================
\begin{block}{The project}
  \begin{itemize}
  \item The CREM laboratory and Parisson have been developing an innovative,
    collaborative and interdisciplinary open-source web-based multimedia
    platform since 2007.
  \item Goal : fit the professional requirements from both sound archivists and
    researchers in ethnomusicology.
  \item Official platform online since 2011 : \emph{Archives sonores du CNRS, Musée de
      l'Homme}: 
  \end{itemize}
\vspace{-0.6cm}
  \begin{center}
    \colorbox{yellow!50} { \hskip3ex \bf \url{http://archives.crem-cnrs.fr} \hskip3ex }
  \end{center}
\end{block}
\begin{block}{The platform}
  \begin{itemize}
  \item \emph{Telemeta} is an \alert{open-source audio web Content
      Management System} (CMS) dedicated to \alert{digital sound
      archives} secure storing, indexing and publishing with database management through (\alert{SQL}) or Oracle backend.
  \item The demonstration presents the features of this platform
    in the context of \alert{\emph{ethnomusicological} research}.
  \item It focuses on the enhance and \alert{collaborative}
    user-experience in accessing audio items and their associated \alert{metadata} and on the possibility for the expert user to further enrich those metadata.\\
    Metadata consist in both \alert{Contextual Ethnomusicological Information} and 
    \alert{Annotations and segmentations} (by human or automatic).
  \item \emph{Telemeta} also provides integrated \alert{audio signal
      processing tools} for automatic analysis of sound items through an external component, \emph{TimeSide}.
  \end{itemize}
\vspace{-0.5cm}
    \begin{center}
      \includegraphics[width=0.3\textwidth]{../poster/img/logo_telemeta_1-1.pdf}\\
      \colorbox{yellow!50}{\textbf{\url{http://telemeta.org/}}}
    \end{center}
  \end{block}
\end{frame}
\begin{frame}\tiny
  \frametitle{TimeSide}
 \begin{block}{TimeSide : open web audio processing framework}
One specificity of the \emph{Telemeta} architecture is to rely on an external component, \emph{TimeSide}, that offers audio player integration together with low and high level audio signal processing capabilities.
\vspace{-0.25cm}
\begin{center}
  \colorbox{yellow!50}{\bf \hskip3ex \url{https://github.com/yomguy/TimeSide/} \hskip3ex  }
\end{center}
\vspace{-0.5cm}
\textbf{Goals}\\
\begin{itemize}
\item \alert{Do} asynchronous and fast audio processing with Python.
\item \alert{Decode} ANY audio or video format into numpy arrays thanks to Gstreamer.
\item \alert{Analyze} audio content with some external audio feature extraction libraries.
\item \alert{Organize}, \alert{serialize} and \alert{save} analysis metadata through various formats.
\item \alert{Draw} various fancy waveforms, spectrograms and other cool graphers.
\item \alert{Transcode} audio data in various media formats and stream them through web apps.
\item \alert{Playback}, \alert{index}, \alert{tag} and \alert{interact} on demand with a smart high-level HTML5 extensible player.
\end{itemize}
\end{block}

 \begin{block}{Audio features extraction}
TimeSide incorporates some state-of-the-art audio feature extraction libraries such as:
\vspace{-0.1cm}
\begin{itemize}
\item \textbf{Aubio:
    \colorbox{yellow!50}{\hskip1ex  \url{http://aubio.org} \hskip1ex }}
\vspace{-0.1cm}
\item \textbf{Yaafe:
    \colorbox{yellow!50}{\hskip1ex \url{http://yaafe.sourceforge.net}\hskip1ex }}
\vspace{-0.1cm}
\item \textbf{Vamp plugins:  
    \colorbox{yellow!50}{\hskip1ex \url{http://www.vamp-plugins.org}\hskip1ex }} 
\end{itemize}
Given the extracted features, every sound item in a given collection can be automatically analyzed. The results of this analysis can be displayed as a support to ethnomusicological studies.
Further works lead by the DIADEMS project will incorporate advanced Music Information Retrieval methods in order to provide \alert{automatic annotation}, \alert{segmentation} and \alert{similarity} analysis.
    \end{block}
\end{frame}
\begin{frame}
  \frametitle{Demonstration}
 
  \begin{itemize}
  \item \textbf{Demo Session \#1}
  \item \alert{Wednesday} October 16th, 14:30 - 16:00
  \item In the Hall
  \end{itemize}
 \begin{center}
    \includegraphics[width=.3\linewidth]{../poster/img/parisson_logo_FINALE_com.pdf}
  \end{center}

  \begin{itemize}
  \item Live demonstration of the \emph{Telemeta} web platform with the CNRS ethnomusicological sound archives
  \item Live demonstration \emph{Timeside} audio processing framework (Python)
  \end{itemize}
\end{frame}
\end{document}
%%% Local Variables: 
%%% mode: latex
%%% TeX-master: t
%%% End: 
